\chapter{Party Drugs}
\index{Party Drugs}

\medskip
\noindent\textit{\textbf{Under review.} This chapter is an AI-assisted educational draft; it is not a substitute for professional medical advice.}

\section*{Definition and Overview}
Party Drugs is a medicine or treatment-related topic. Its safe use depends on the indication, dose, interactions, contraindications, and monitoring required for the individual.

\section*{Clinical Features}
Possible effects and adverse effects vary with the medicine and may include gastrointestinal, neurological, allergic, cardiovascular, or organ-specific problems.

\section*{When to Seek Medical Care}
Seek medical review for new, persistent, worsening, or function-limiting symptoms. Contact your local emergency services for severe breathing difficulty, collapse, sudden neurological change, severe chest or abdominal pain, or any rapidly worsening illness.

\section*{Diagnosis and Investigations}
Review the indication, medicines, allergies, kidney and liver function, pregnancy status, and relevant monitoring tests.

\section*{Management}
Use only as directed by a qualified clinician or pharmacist. Do not share medicines, change doses, or stop long-term treatment without advice.

\section*{Complications}
Complications depend on the underlying cause and may include progression, effects on other organs, treatment adverse effects, or reduced quality of life.

\section*{Prognosis and Self-Care}
Outcome depends on the cause, severity, complications, and response to treatment. Follow the care plan, keep appointments, avoid known triggers, and seek help if symptoms worsen.
\section*{Safety-netting}
Seek urgent help for severe or rapidly worsening symptoms, collapse, breathing difficulty, severe pain, confusion, dehydration, uncontrolled bleeding, or any immediate danger.
